% Sección: formulación del problema como MDP.
\section{Modelo: formulación como MDP}
\label{sec:mdp}

Modelamos Tetris como un proceso de decisión de Markov (MDP) episódico
$\mathcal{M} = (\mathcal{S}, \mathcal{A}, P, R, \gamma)$. El carácter \emph{episódico} proviene de
que toda partida termina en un estado absorbente (\emph{top-out}) tras un número
finito de pasos, lo que da un retorno bien definido sin necesidad de descuento.

\subsection{Elementos}
\begin{itemize}[itemsep=2pt]
    \item \textbf{Estado} $s = (b, p)$: tablero $b \in \{0,1\}^{H\times W}$
          ($H=20$, $W=10$) y pieza actual $p \in \{1,\dots,7\}$. Un estado
          alcanzable nunca contiene una fila completa, pues se elimina al fijar.
    \item \textbf{Acciones} $\mathcal{A}(s) = \{(r,c)\}$: una \emph{colocación} es
          una rotación $r$ y una columna de anclaje $c$ factibles para la pieza
          actual (Sección~\ref{subsec:acciones}).
    \item \textbf{Transición} $P(s' \mid s, a)$: factoriza en una parte
          determinista y una aleatoria (Sección~\ref{subsec:transicion}).
    \item \textbf{Recompensa} $R(s, a) = L(s, a)$: líneas eliminadas por la
          colocación (Sección~\ref{subsec:recompensa}).
    \item \textbf{Descuento} $\gamma = 1$ (Sección~\ref{subsec:decisiones}).
\end{itemize}

\subsection{Terminología}
Un \emph{tetrominó} es una pieza de cuatro celdas; empleamos los siete estándar
($I, O, T, S, Z, J, L$), cada uno con entre una y cuatro \emph{rotaciones}. El
\emph{hard-drop} deja caer la pieza en línea recta hasta que apoya y la
\emph{fija} (\emph{lock}); el \emph{soft-drop} es una caída controlada que todavía
permite maniobrar. Al fijarse, toda fila completa se elimina (\emph{line clear})
y las superiores descienden. El \emph{afterstate} $\sigma$ es el tablero
resultante tras fijar y limpiar, antes de revelar la pieza siguiente
(\emph{spawn}). La partida termina en \emph{top-out} cuando la nueva pieza no
admite ninguna colocación. Llamamos \emph{hueco} a una celda vacía con alguna
celda ocupada por encima, \emph{pozo} a una columna estrecha y profunda, y
\emph{altura} de la columna $j$ a $h_j = H - \min\{\, i : b_{i,j} = 1 \,\}$.

\subsection{Acciones por colocación: finitud y cota}
\label{subsec:acciones}
Cada acción es un par $(r, c)$. El número de rotaciones distintas de cualquier
tetrominó es a lo sumo $4$. Fijada una rotación de ancho $w_r$ con
$1 \le w_r \le 4$, las columnas de anclaje válidas son $W - w_r + 1 \le W$. En
consecuencia,
\[
    |\mathcal{A}(s)| \;\le\; \sum_{r} (W - w_r + 1) \;\le\; 4W,
\]
y $\mathcal{A}(s)$ es finito por ser un producto de conjuntos finitos (rotaciones
y columnas). El conteo real es menor, ya que se descartan las colocaciones que
desbordan la pila. La formulación por colocación reduce el control a \emph{una
    decisión por pieza}.

\subsection{Transición y mapa de colocación}
\label{subsec:transicion}
El estado siguiente tiene dos componentes, $s' = (\sigma', p')$ (el tablero y la
pieza siguiente), así que la transición es una distribución \emph{conjunta}. La
descomponemos con la regla de la cadena:
\begin{equation}
    P\big( (\sigma', p') \mid s, a \big)
    = \mathbb{P}(\sigma' \mid s, a)\;\cdot\;\mathbb{P}(p' \mid \sigma', s, a).
    \label{eq:cadena}
\end{equation}
El \emph{mapa de colocación} $\sigma = f(s, a)$ es determinista: dado $(s, a)$, el
hard-drop más la limpieza de líneas producen un \emph{único} tablero. Una
distribución con toda su masa en un solo valor es una delta, que escribimos con la
función indicadora
\begin{equation}
    \mathbf{1}\{C\} =
    \begin{cases}
        1 & \text{si la condición } C \text{ es verdadera}, \\[2pt]
        0 & \text{en caso contrario},
    \end{cases}
\end{equation}
de modo que el primer factor de \eqref{eq:cadena} es
\begin{equation}
    \mathbb{P}(\sigma' \mid s, a) = \mathbf{1}\{\sigma' = f(s, a)\}.
    \label{eq:delta}
\end{equation}
El segundo factor es la pieza siguiente, que se sortea de forma independiente del
tablero, de la acción y de la historia (es exógena), y la modelamos uniforme sobre
las siete:
\begin{equation}
    \mathbb{P}(p' \mid \sigma', s, a) = \mathbb{P}(p') = \tfrac{1}{7},
    \qquad p' \sim \mathrm{Unif}\{1, \dots, 7\}.
    \label{eq:pieza}
\end{equation}
Sustituyendo \eqref{eq:delta} y \eqref{eq:pieza} en \eqref{eq:cadena} se obtiene la
transición:
\begin{equation}
    P\big( (\sigma', p') \mid s, a \big)
    = \mathbf{1}\{\sigma' = f(s, a)\}\;\cdot\;\tfrac{1}{7}.
    \label{eq:transicion}
\end{equation}
Es decir, desde $(s, a)$ solo son alcanzables los siete sucesores
$(\sigma, 1), \dots, (\sigma, 7)$ y equiprobables; para cualquier $\sigma' \neq f(s, a)$ la probabilidad es
nula. Elegir el generador independiente (i.i.d.) preserva la propiedad de Markov
sin ampliar el estado: un aleatorizador por bolsas (\emph{7-bag}) obligaría a
recordar las piezas ya emitidas. Esta factorización es la que habilita trabajar con el afterstate $\sigma$
(Sección~\ref{sec:valor}).

\subsection{Recompensa y su rango}
\label{subsec:recompensa}
$L(s, a)$ es el número de filas completas eliminadas por la colocación. Como un
tetrominó ocupa cuatro celdas repartidas en a lo sumo cuatro filas, y una fila
solo puede completarse si la pieza aporta alguna celda a ella, se eliminan a lo
sumo cuatro filas. El mínimo $0$ (ninguna fila completada) y el máximo $4$ (una
$I$ vertical que cierra cuatro filas a la vez, un \emph{Tetris}) son alcanzables,
de modo que
\[
    L(s, a) \in \{0, 1, 2, 3, 4\}.
\]

\subsection{Decisiones de modelado}
\label{subsec:decisiones}
\paragraph{Recompensa.} Tomamos $R = L$ porque el objetivo del problema es
maximizar el total de líneas de la partida: es una señal acotada y directamente
alineada con la métrica de evaluación.

\paragraph{Descuento $\gamma = 1$.} Toda línea vale igual sin importar cuándo se
elimina; un $\gamma < 1$ introduciría un sesgo temporal (preferir limpiar antes)
ajeno al objetivo, que además penalizaría acumular para un \emph{Tetris}. Al ser
la tarea episódica y terminar casi seguramente, la suma no descontada es finita;
y el optimizador (CEM) maximiza $J(w) = \mathbb{E}[\sum_t R_t]$ de forma directa,
sin \emph{bootstrapping}, por lo que $\gamma = 1$ equivale a ``sumar las líneas''.
En evaluación se acota la longitud del episodio para mantener $J$ finito.

\paragraph{Abstracción del tiempo.} El modelo prescinde de la tasa de refresco
(\emph{frame rate}) y del tiempo real: un paso del MDP equivale a \emph{una pieza
    colocada}, y la dinámica intra-pieza (gravedad, soft-drop, desplazamientos
laterales) se colapsa en la elección de la posición final más el hard-drop. Se
asume, pues, una \emph{decisión inmediata}: lo relevante para la calidad de la
política es \emph{dónde} reposa la pieza, no la trayectoria ni el cronometraje
para llevarla allí. Como implicación directa, \textbf{no se modelan las maniobras
    que requieren rotar o desplazar la pieza durante la caída} (\emph{tucks},
\emph{T-spins}): $\mathcal{A}(s)$ enumera únicamente las caídas rectas alcanzables
por hard-drop. Tampoco se modelan la aceleración de la gravedad por nivel ni el
puntaje por soft-drop. El bucle de frames solo sería relevante para una
demostración en vivo, donde la colocación elegida se traduciría a una secuencia
de entradas de control.
